% Source: https://artofproblemsolving.com/community/c3h2428579_imperfect_squares
% Stable identity: aops-topic-2428579-q01
\textbf{(a)} Suppose, for contradiction, that $\sqrt n=\dfrac pq$ with positive integers $p,q$ that have no common factor greater than $1$. Squaring gives
\[p^2=nq^2.\]
If $q>1$, let $r$ be a prime factor of $q$. Then $r$ divides $nq^2=p^2$, and since $r$ is prime it divides $p$. So $r$ is a common factor of $p$ and $q$, a contradiction. Hence $q=1$ and $n=p^2$ is a perfect square, contradicting the assumption. Therefore $\sqrt n$ is irrational. $\blacksquare$

\textbf{(b)} Every terminating or eventually repeating decimal is rational:
\begin{itemize}
\item a terminating decimal with $k$ digits after the point equals $\dfrac{N}{10^k}$ for an integer $N$;
\item if $x$ has a block of $m$ digits that repeats forever after the first $k$ digits after the point, then $10^{k+m}x$ and $10^kx$ have the same digits after the point, so $10^{k+m}x-10^kx=N$ is an integer and $x=\dfrac{N}{10^k\left(10^m-1\right)}$.
\end{itemize}
Since $\sqrt n$ is irrational by (a), its decimal expansion can be neither terminating nor eventually repeating. For example, $\sqrt2=1.41421356\ldots$ never settles into a repeating block.
